\section{Compact-Fiber Carrier and Electromagnetic Bridge}

\subsection{Compact-fiber carrier}

The compact-fiber carrier is
\[
F=Z_4\times Z_4\times Z_8.
\]
This carrier and its sector structure are established in the foundational compact-fiber work and used in subsequent atomic, representational, and radiation-sector applications \cite{Wells2026PaperI,Wells2026PaperII,Wells2026PaperIII,Wells2026Radiation}.

Its cardinality is
\[
|F|=4\cdot4\cdot8=128.
\]

The two \(Z_4\) factors supply the paired magnetic sector. The \(Z_8\) factor supplies the electric-cover sector. The corresponding sector capacities are
\[
N_m=4,
\qquad
N_e=8.
\]

The relation
\[
N_m\mid N_e,
\qquad
8=2\cdot4,
\]
gives the canonical electric-cover / magnetic-base quotient
\[
\pi:Z_8\to Z_4.
\]
This quotient is not independent bridge data. It is downstream of the compact-fiber sector capacities. In the alpha derivation, \(\pi\) becomes relevant when paired magnetic feedback is read through the electric-cover structure.

\subsection{Material/radiative bridge}

The electromagnetic coupling is treated as a carrier-level bridge normalization between material closure and radiative/vibrational placement.

Material closure is represented by
\[
F=Z_4\times Z_4\times Z_8.
\]
Radiative/vibrational placement is represented by a bridge-placement carrier
\[
V_{\rm rad}^{\rm bridge}.
\]
The electromagnetic bridge is therefore the material/radiative interface:
\[
F
\longleftrightarrow
V_{\rm rad}^{\rm bridge}.
\]

The coupling derived here is denoted
\[
\alpha_F.
\]
Its inverse is
\[
x=\alpha_F^{-1}.
\]
The quantity \(x\) is the effective compact-fiber electromagnetic bridge capacity.

\subsection{Carrier-level universality}

Within the compact-fiber framework, the electromagnetic bridge is a carrier-to-carrier material/radiative interface. Its coupling is therefore universal at the bridge level. For admissible element labels
\[
f\in F,
\qquad
v\in V_{\rm rad}^{\rm bridge},
\]
the bridge coupling satisfies
\[
\alpha_F(f,v)=\alpha_F.
\]

Thus the electromagnetic coupling is not an element-level quantity whose value depends on a selected compact-fiber state and a selected radiative-placement state. It is a scalar per-channel bridge share determined by carrier-level access capacity.

This carrier-level universality is the reason the alpha bridge is not normalized over the pair-configuration space
\[
F\times V_{\rm rad}^{\rm bridge}.
\]
A pair space would be appropriate for an observable whose elementary alternatives are ordered pairs \((f,v)\). The alpha bridge has a different role: it is the scalar normalization of the material/radiative interface itself.

The relevant structural inputs are therefore carrier-level invariants:
\[
|F|,
\qquad
|V_{\rm rad}^{\rm bridge}|,
\qquad
N_m,
\qquad
N_e.
\]
For the compact fiber,
\[
|F|=128,
\qquad
N_m=4,
\qquad
N_e=8.
\]
The remaining carrier input is the radiative/vibrational bridge-placement count
\[
|V_{\rm rad}^{\rm bridge}|,
\]
which is derived below.

\subsection{\texorpdfstring{Why the bridge is not counted as \(F\times V_{\rm rad}^{\rm bridge}\)}{Why the bridge is not counted as F x Vrad bridge}}

A Cartesian-product count
\[
|F\times V_{\rm rad}^{\rm bridge}|
=
|F|\,|V_{\rm rad}^{\rm bridge}|
\]
would be appropriate if an electromagnetic bridge state required simultaneous specification of one material element and one radiative-placement element as a joint state label.

That is not the role of the electromagnetic bridge in this derivation.

The bridge coupling is not an element-pair table. It is a scalar carrier-level bridge normalization. The material carrier and the radiative/vibrational bridge-placement carrier supply mechanism-distinct access to the bridge. They are independently sufficient carrier-level access structures, not simultaneously required labels of a joint microscopic state.

Therefore the base bridge count uses disjoint-union access,
\[
F\sqcup V_{\rm rad}^{\rm bridge},
\]
not product-state access,
\[
F\times V_{\rm rad}^{\rm bridge}.
\]
Thus the base count is
\[
B=|F|+|V_{\rm rad}^{\rm bridge}|.
\]

This distinction is load-bearing. It separates the compact-fiber alpha derivation from a pair-counting argument.

\subsection{\texorpdfstring{Structural role of \(x=\alpha_F^{-1}\)}{Structural role of x = alphaF inverse}}

The quantity
\[
x=\alpha_F^{-1}
\]
is not a primitive compact-fiber cardinality. It is the self-consistent effective bridge capacity obtained after the base carrier count is corrected by cover-mediated feedback.

The derivation proceeds in two stages. First, the base bridge count
\[
B
\]
is obtained from carrier-level material and radiative/vibrational bridge-placement access. Second, the net cover-mediated feedback burden
\[
K
\]
is shared over the final effective bridge capacity
\[
x.
\]

This gives the first-order fixed point
\[
x=B+\frac{K}{x}.
\]
A closed-return correction then refines this to
\[
x
=
B+\frac{K}{x}
-
\frac{K N_m}{(N_e-1)x^3}.
\]

The positive root of the resulting quartic is the compact-fiber structural prediction for
\[
\alpha_F^{-1}.
\]

\subsection{Structural output}

The output of this derivation is the compact-fiber electromagnetic bridge coupling
\[
\alpha_F.
\]
The later identification
\[
\alpha_F=\alpha_{\rm low-energy}
\]
is made by physical role, not by using the measured fine-structure constant as an input.